\subsection{Rotation invariance and nontriviality: what the lattice inputs
cannot give}\label{rN5:sec}

Section~\ref{ssec:limitE} reduces the Euclidean invariance of a continuum
limit to the one asymptotic symmetry (W$_{\mathrm s}$): by
Corollary~\ref{cor:limitEexact} the moments converge to an
$\mathrm O(4)$-invariant family if and only if (M$_{\mathrm s}$),
(W$_{\mathrm s}$) and (U$_{\mathrm s}$) hold, and
Theorem~\ref{thm:limitErot} derives (W$_{\mathrm s}$) from the Symanzik-type
inputs (SB$_n$), (RI$_n$), which are not proved for the Wilson measure. This
subsection answers two questions about this reduction.

First, whether (W$_{\mathrm s}$) could be a \emph{consequence} of the
structural properties that the lattice inputs do provide: invariance under
translations and $B_4$, reflection positivity for both reflection types in
all four coordinate directions, a transfer matrix in every direction with a
gap uniform in physical units, exponential clustering, uniqueness of the
limit, convergence along the whole sequence, the moment bounds (M),
(M$_{\mathrm s}$) and (G$_{\mathrm s}$), the Osterwalder--Schrader
temperedness and linear growth conditions, and real analyticity of the
limit off the coincidence set. The Gaussian fields of
Proposition~\ref{prop:osobstruction}(a) (Remark~\ref{rem:limitEscope}) and
of Proposition~\ref{prop:limitEafexample} already have $G_S=B_4$, but their
kernels are functions of $\|v\|_1$, not differentiable on the hyperplanes
$v_\mu=0$, and the cube condition is not verified for them; so one could
still hope that analyticity of the limit, or the cube condition itself,
forces $G_S=\mathrm O(4)$. Theorem~\ref{rN5:nogo} shows that it does not.
The field is the average $\phi=\frac12\sum_\mu\phi^{(\mu)}$ of the
components of a four-component Gaussian field with the nearest-neighbour
gradient action, the component $\phi^{(\mu)}$ having speed of light
$\lambda\ne1$ in the direction $\mu$ and $1$ in the others; $B_4$ acts on
the sites and permutes the components. It satisfies the cube condition on
tori uniformly at fixed physical cube size (Lemma~\ref{rN5:bp}), has all
the properties listed above, and its limit, a free field with four speeds
of light, has a real analytic kernel and $G_S=B_4$ exactly. Hence no
theorem of the form ``reflection positivity, cubic symmetry, cube
condition, mass gap, clustering, regularity $\Rightarrow$ $\mathrm O(4)$''
exists, and the restoration of rotation invariance must come from a
dynamical input specific to the Wilson measure, of the type (RI$_n$) or
(W$_{\mathrm s}$) itself. (The example is not a lattice gauge theory; what
it shows is that a proof of (W$_{\mathrm s}$) cannot use only these
properties.)

Second, what the gap and the cube estimate contribute to nontriviality. They
exclude the Schwinger functions of Wick polynomials of \emph{massless} free
fields, in particular of the free Maxwell field of $\U(1)^{\dim G}$
(Proposition~\ref{rN5:nomassless}); they do not exclude massive free
fields, which satisfy the cube estimate uniformly
(Theorems~\ref{thm:ffscalar} and~\ref{thm:ffproca}); Gaussian gauge
fields have a perimeter law (Lemma~\ref{rN5:gaussloop}); and a bounded
nonzero physical string tension along the tuning excludes every Gaussian
gauge field with summable covariance (Theorem~\ref{rN5:stgauss}), but is
not implied by the cube estimate and the gap and is not proved here.
Throughout, as in Proposition~\ref{prop:limitEafexample}, the hypotheses of
Section~\ref{ssec:limitE} on the states $\omega_k$ are read with
``function of the field at the sites'' in place of ``gauge-invariant
function of the links'', with $r=0$; by Remark~\ref{rem:limitEscope} the
cube condition enters that section only through
Lemma~\ref{lem:limitElattice}(iii).

\subsubsection{The anisotropic lattice Green function}

Throughout this part $T^d=[-\pi,\pi]^d$ carries the measure $d^dp/(2\pi)^d$.
Let $c=(c_0,\dots,c_3)$ with $c_\nu>0$, $c_-=\min_\nu c_\nu$, $c_+=\max_\nu
c_\nu$, let $M\ge0$ and, for $z\in\Z^4$,
\begin{equation}\label{rN5:eq:green}
C_{M,c}(z)=\int_{T^4}\frac{e^{ip\cdot z}}{M^2+\sum_\nu c_\nu\,2(1-\cos
p_\nu)}\,\frac{d^4p}{(2\pi)^4}.
\end{equation}
For $M>0$ this is the covariance of the centred Gaussian field $\xi$ on
$\Z^4$ with Gibbs action $\frac12\sum_x\bigl[M^2\xi_x^2+\sum_\nu c_\nu(\xi_{x+e_\nu}
-\xi_x)^2\bigr]$: the precision operator $M^2-\sum_\nu c_\nu\Delta_\nu$,
$\Delta_\nu$ the one-dimensional lattice Laplacian in direction $\nu$, has
the Fourier symbol in the denominator of~\eqref{rN5:eq:green}, which is
bounded below by $M^2>0$, so the field exists on all of $\Z^4$ and is the
limit of the Gibbs measures on tori. We write $\vec p=(p_1,p_2,p_3)$ and
$\vec z$ for the spatial parts and $\omega(\vec p)^2=M^2+\sum_{i=1}^3c_i\,2(1-\cos p_i)$.

\begin{lemma}\label{rN5:green}
Let $M\in[0,1]$. Define $\varepsilon(\vec p)\ge0$ by $\cosh\varepsilon(\vec
p)=1+\omega(\vec p)^2/(2c_0)$.

\textup{(i)} For every $z\in\Z^4$ with $\omega(\vec p)>0$ for almost every
$\vec p$ \textup{(}always when $M>0$\textup{)},
\begin{equation}\label{rN5:eq:3d}
C_{M,c}(z)=\int_{T^3}e^{i\vec p\cdot\vec z}\,\frac{e^{-\varepsilon(\vec
p)|z_0|}}{2c_0\sinh\varepsilon(\vec p)}\,\frac{d^3p}{(2\pi)^3}.
\end{equation}

\textup{(ii)} $\sinh\varepsilon\ge\omega/\sqrt{c_0}$ and $\varepsilon\ge
\kappa\,\omega$ with $\kappa=(c_0+3c_++\tfrac14)^{-1/2}$, and $\omega(\vec
p)^2\ge M^2+(4c_-/\pi^2)|\vec p|^2$.

\textup{(iii)} There are $K_0,K_2,K_3<\infty$ depending only on $c_\pm$
such that for all $M\in[0,1]$ and $z\in\Z^4$: $|C_{M,c}(z)|\le C_{M,c}(0)\le
K_0$; $|C_{M,c}(z)|\le K_2/|z|_\infty^2$ for $z\ne0$; and
$|C_{M,c}(z)|\le K_3e^{-\varepsilon(0)|z_0|}$ with
\[
\varepsilon(0)=\operatorname{arccosh}\Bigl(1+\frac{M^2}{2c_0}\Bigr)\ge\frac
M{\sqrt{c_0}}\Bigl(1+\frac{M^2}{4c_0}\Bigr)^{-1/2} .
\]
The same bounds hold with any coordinate direction in the role of the time
direction $0$ \textup{(}with $c_0$ replaced by that $c_\nu$ in
$\varepsilon(0)$\textup{)}, and the first two hold for every $M\ge0$.

\textup{(iv)} Let $m>0$, $a_k\downarrow0$, $M_k=ma_k\le1$, and $W_c(y)=
\int_{\R^3}e^{i\vec q\cdot\vec y}\,e^{-\Omega(\vec q)|y_0|}(2c_0\Omega(\vec
q))^{-1}\,d^3q/(2\pi)^3$ for $y_0\ne0$, with $\Omega(\vec q)=\bigl((m^2+
\sum_ic_iq_i^2)/c_0\bigr)^{1/2}$. If $z_k\in\Z^4$ and $a_kz_k\to y$ with
$y_0\ne0$, then $a_k^{-2}C_{M_k,c}(z_k)\to W_c(y)$, uniformly for $y$ in
compact subsets of $\{y_0\ne0\}$.
\end{lemma}

\begin{proof}
(i) For $\varepsilon>0$ and real $\theta$, summing two geometric series,
\[
\sum_{n\in\Z}e^{-\varepsilon|n|}e^{-in\theta}=\frac{1-e^{-2\varepsilon}}{1-2e^
{-\varepsilon}\cos\theta+e^{-2\varepsilon}}=\frac{\sinh\varepsilon}{\cosh
\varepsilon-\cos\theta},
\]
so the Fourier coefficients of $\theta\mapsto(\cosh\varepsilon-\cos\theta)^{-1}$
are $e^{-\varepsilon|n|}/\sinh\varepsilon$. The denominator
of~\eqref{rN5:eq:green} is $\omega(\vec p)^2+2c_0(1-\cos p_0)=2c_0(\cosh
\varepsilon(\vec p)-\cos p_0)$ by the definition of $\varepsilon$. Integrating
first over $p_0$, for each $\vec p$ with $\omega(\vec p)>0$,
\[
\int_{-\pi}^\pi\frac{e^{ip_0z_0}}{2c_0(\cosh\varepsilon-\cos p_0)}\frac
{dp_0}{2\pi}=\frac{e^{-\varepsilon|z_0|}}{2c_0\sinh\varepsilon},
\]
and Fubini applies because the integrand of~\eqref{rN5:eq:green} is
integrable on $T^4$ (the denominator is bounded below by
$(4c_-/\pi^2)|p|^2$ by the bound of (ii), and $|p|^{-2}$ is integrable in
four dimensions).

(ii) $1-\cos\theta\ge2\theta^2/\pi^2$ on $[-\pi,\pi]$ gives the bound on
$\omega^2$. From $\cosh\varepsilon-1=\omega^2/(2c_0)=:u$,
$\sinh^2\varepsilon=(\cosh\varepsilon-1)(\cosh\varepsilon+1)\ge2u=\omega^2/c_0$.
Further $\varepsilon=2\operatorname{arcsinh}\sqrt{u/2}$ and
$\operatorname{arcsinh}s=\int_0^s(1+t^2)^{-1/2}dt\ge s(1+s^2)^{-1/2}$, so
$\varepsilon\ge\sqrt{2u}\,(1+u/2)^{-1/2}=\omega\,(c_0+\omega^2/4)^{-1/2}$;
since $\omega^2\le M^2+12c_+\le1+12c_+$ this is at least $\kappa\omega$.

(iii) By~\eqref{rN5:eq:3d} and (ii), with $\varepsilon\ge\kappa\omega\ge\kappa
(2\sqrt{c_-}/\pi)|\vec p|=:\kappa'|\vec p|$ and $\sinh\varepsilon\ge\omega/
\sqrt{c_0}\ge(2\sqrt{c_-}/\pi)|\vec p|/\sqrt{c_0}$,
\[
\begin{split}
|C_{M,c}(z)|&\le\frac{\pi}{4\sqrt{c_0c_-}}\int_{\R^3}\frac{e^{-\kappa'|\vec
p||z_0|}}{|\vec p|}\frac{d^3p}{(2\pi)^3}\\
&=\frac{1}{8\pi\sqrt{c_0c_-}}\int_0^\infty\rho e^{-\kappa'\rho|z_0|}d\rho
=\frac{1}{8\pi\sqrt{c_0c_-}\,\kappa'^2z_0^2}
\end{split}
\]
for $z_0\ne0$. Since~\eqref{rN5:eq:green} is symmetric under permutations of
the coordinates together with the $c_\nu$, the same argument with the
direction $\nu$ of $|z|_\infty=|z_\nu|$ as time gives $|C_{M,c}(z)|\le
K_2/|z|_\infty^2$ with $K_2$ depending only on $c_\pm$. The bound $C_{M,c}(0)\le
K_0$ follows from~\eqref{rN5:eq:green} and $\omega^2+2c_0(1-\cos p_0)\ge(4c_-/
\pi^2)|p|^2$, and $|C_{M,c}(z)|\le C_{M,c}(0)$ because $C_{M,c}$ is the Fourier
transform of a positive integrable function. For the exponential bound,
$\varepsilon(\vec p)\ge\varepsilon(0)$ since $\omega(\vec p)\ge\omega(0)=M$,
so~\eqref{rN5:eq:3d} gives $|C_{M,c}(z)|\le e^{-\varepsilon(0)|z_0|}\int_{T^3}
(2c_0\sinh\varepsilon)^{-1}\le e^{-\varepsilon(0)|z_0|}\int_{T^3}\frac{\pi}
{4\sqrt{c_0c_-}|\vec p|}\frac{d^3p}{(2\pi)^3}=:K_3e^{-\varepsilon(0)|z_0|}$;
the lower bound for $\varepsilon(0)$ is the inequality of (ii) at $\omega=M$.
For $M>1$ the absolute value of the integrand of~\eqref{rN5:eq:3d} is at
most its value at $M=1$, since $\varepsilon(\vec p)$ increases with $M$,
so the bounds derived from it at $M=1$ remain valid; and $C_{M,c}(0)\le
C_{1,c}(0)$ by~\eqref{rN5:eq:green}.

(iv) Substitute $\vec p=a_k\vec q$ in~\eqref{rN5:eq:3d}:
\[
a_k^{-2}C_{M_k,c}(z_k)=\int_{|q_i|\le\pi/a_k}e^{i\vec q\cdot a_k\vec z_k}\,
\frac{a_k\,e^{-(\varepsilon(a_k\vec q)/a_k)\,a_k|z_{k,0}|}}{2c_0\sinh
\varepsilon(a_k\vec q)}\,\frac{d^3q}{(2\pi)^3}.
\]
As $k\to\infty$, $\omega(a_k\vec q)^2/a_k^2=m^2+\sum_ic_i\,2(1-\cos a_kq_i)/a_k^2
\to m^2+\sum_ic_iq_i^2$, hence $u_k=\omega(a_k\vec q)^2/(2c_0)\to0$ and, from
$\varepsilon=2\operatorname{arcsinh}\sqrt{u/2}$ and $\sinh\varepsilon=\sqrt
{2u(1+u/2)}$, both $\varepsilon(a_k\vec q)/a_k$ and $\sinh\varepsilon(a_k\vec
q)/a_k$ tend to $\Omega(\vec q)$; and $a_k|z_{k,0}|\to|y_0|$. So the integrand
converges pointwise to that of $W_c(y)$, uniformly for $y$ in a compact
$\mathcal K\subset\{y_0\ne0\}$ (the dependence on $y$ enters only through
$a_kz_k$, and $|a_kz_k-y|\le a_k$ may be assumed). For the domination let
$\delta=\min_{\mathcal K}|y_0|>0$ and $k$ so large that $a_k|z_{k,0}|\ge
\delta/2$ for all $y\in\mathcal K$. By (ii),
\[
\frac{a_k}{2c_0\sinh\varepsilon(a_k\vec q)}\le\frac{a_k\sqrt{c_0}}{2c_0\,
\omega(a_k\vec q)}\le\frac\pi{4\sqrt{c_0c_-}\,|\vec q|}
\]
and
$e^{-\varepsilon|z_{k,0}|}\le e^{-\kappa'a_k|\vec q||z_{k,0}|}\le e^{-\kappa'
\delta|\vec q|/2}$, so the integrands are dominated by $\pi e^{-\kappa'\delta
|\vec q|/2}/(4\sqrt{c_0c_-}|\vec q|)\in L^1(\R^3)$ uniformly in $y\in\mathcal
K$. Dominated convergence applied to $\sup_{y\in\mathcal K}$ of the difference
of the integrands gives the uniform convergence.
\end{proof}

\subsubsection{Reflection positivity and the gap of a Gaussian lattice field}

Let $\xi=(\xi_x)_{x\in\Z^4}$ be a centred Gaussian field with covariance
$C(x,y)$ and let $\mathbb E$ denote its expectation. For a map
$\vartheta\colon\Z^4\to\Z^4$ with $\vartheta^2=1$ and $C(\vartheta x,\vartheta
y)=C(x,y)$, the law of $\xi$ is invariant under $\xi\mapsto\xi\circ\vartheta$,
and for a function $F$ of the field we put $\vartheta(F)=F\circ\vartheta$,
that is $\vartheta(F)(\xi)=F(\xi\circ\vartheta)$. The two reflections of
Section~\ref{ssec:limitE} in direction $\mu$ are the \emph{site reflection}
$\theta_\mu$, $(\theta_\mu x)_\mu=-x_\mu$, with positive side $\Lambda^+_\mu=
\{x_\mu\ge0\}$, and the \emph{link reflection} $\vartheta_\mu$, $(\vartheta_\mu
x)_\mu=1-x_\mu$, with positive side $\Lambda^{+}_{\mu,1}=\{x_\mu\ge1\}$; both
fix the other coordinates.

\begin{lemma}[reflection positivity of Gaussian fields]\label{rN5:gaussrp}
Let $\vartheta$ be as above, let $\Lambda^+\subset\Z^4$, and suppose that
the kernel $(x,y)\mapsto
C(\vartheta x,y)$ is positive semidefinite on $\Lambda^+$, that is
$\sum_{x,y\in\Lambda^+}\bar v_xv_yC(\vartheta x,y)\ge0$ for all finitely supported
$v\colon\Lambda^+\to\C$. Then $\mathbb E[\vartheta(\bar F)F]\ge0$ for every
$F\in L^2$ measurable with respect to $(\xi_x)_{x\in\Lambda^+}$.
\end{lemma}

\begin{proof}
For finitely supported real $v,w$ on $\Lambda^+$ put $\xi(v)=\sum_xv_x\xi_x$.
Since $\xi(v)\circ\vartheta=\xi(\vartheta v)$ with $(\vartheta v)_x=v_{\vartheta
x}$ and $\xi(\vartheta v)+\xi(w)$ is centred Gaussian,
\[
\begin{split}
\mathbb E\bigl[\vartheta(\overline{e^{i\xi(v)}})e^{i\xi(w)}\bigr]&=\mathbb E\,
e^{i(\xi(w)-\xi(\vartheta v))}\\
&=e^{-\frac12\ip v{Cv}}e^{-\frac12\ip w{Cw}}\,e^{\ip{\vartheta v}{Cw}},
\end{split}
\]
using $\ip{\vartheta v}{C\vartheta v}=\ip v{Cv}$. The kernel $(v,w)\mapsto\ip
{\vartheta v}{Cw}=\sum_{x,y}v_xw_yC(\vartheta x,y)$ is positive semidefinite
by hypothesis, hence so is its exponential (the Schur product theorem applied
to the partial sums of the exponential series) and so is the product with the
rank-one kernel $e^{-\frac12\ip v{Cv}}e^{-\frac12\ip w{Cw}}$. Thus
$\mathbb E[\vartheta(\bar F)F]\ge0$ for every $F$ in the linear span
$\mathcal T$ of the functions $e^{i\xi(v)}$, $v$ real and finitely supported
on $\Lambda^+$. The sesquilinear form $(F,G)\mapsto\mathbb E[\vartheta(\bar
F)G]$ is bounded on $L^2$, $|\mathbb E[\vartheta(\bar F)G]|\le\|F\|_2\|G\|_2$,
by the Cauchy--Schwarz inequality and the $\vartheta$-invariance of the law.
It remains to see that $\mathcal T$ is dense in $L^2(\sigma(\xi_x:x\in\Lambda^+))$.
Functions of finitely many coordinates $\xi_x$, $x\in\Lambda^+$, are dense
by martingale convergence, so it suffices to show that for a finite
$\Lambda\subset\Lambda^+$ the trigonometric polynomials are dense in
$L^2(\R^\Lambda,\gamma)$, $\gamma$ the (possibly degenerate) Gaussian law of
$(\xi_x)_{x\in\Lambda}$. If $g\in L^2(\gamma)$ is orthogonal to all $e^{i\ip v\cdot}$,
the finite complex measure $g\gamma$ has vanishing Fourier transform, so
$g\gamma=0$ and $g=0$ $\gamma$-a.e.
\end{proof}

The Osterwalder--Schrader reconstruction of Appendix~\ref{app:transfer}
applies verbatim to a translation-invariant Gaussian field on $\Z^4$ that is
reflection positive for $\theta_0$ and $\vartheta_0$: the Hilbert space $\mathcal H$ is the
completion of the functions of $(\xi_x)_{x\in\Lambda^+_{0,1}}$ modulo the null
space of $\ip FG=\mathbb E[\vartheta_0(\bar F)G]$, with $\Omega$ the class of
$1$, and the transfer matrix $\mathbb T$ is induced by the unit translation
$\tau_1$ in direction $0$, $0\le\mathbb T\le1$. The proof of
Proposition~\ref{prop:gapdecay} uses only this structure, so its criterion
holds for such fields; we use it with the class of polynomials in finitely
many $\xi_x$, $x\in\Lambda^+_{0,1}$, which is dense in $L^2$ on each finite
set of sites $\Lambda$: the trigonometric polynomials are dense in
$L^2(\R^\Lambda,\gamma)$ (proof of Lemma~\ref{rN5:gaussrp}), and each $e^{i
\ip v\cdot}$ is the $L^2(\gamma)$-limit of its Taylor polynomials, since these
are dominated by $e^{|v||\cdot|}\in L^2(\gamma)$; see also~\cite{Jan97}.

\begin{lemma}[gap of a Gaussian lattice field]\label{rN5:gaussgap}
Let $\xi$ be a translation-invariant centred Gaussian field on $\Z^4$ with
covariance $C(x-y)$, reflection positive for $\theta_0$ and $\vartheta_0$,
and suppose that
$|C(z)|\le Ke^{-\epsilon|z_0|}$ for all $z$. Then the spectrum of $\mathbb
T$ on $\Omega^\perp$ lies in $[0,e^{-\epsilon}]$.
\end{lemma}

\begin{proof}
Let $F=\prod_{i=1}^n\xi_{x_i}$ and $F'=\prod_{j=1}^{n'}\xi_{y_j}$ with
$x_i,y_j\in\Lambda^+_{0,1}$, and $t\ge0$. By Isserlis' formula~\cite{Iss18}
(Wick's theorem), $\mathbb E[\vartheta_0(F)\tau_tF']$ is the sum over the
pairings of the $n+n'$ points $\vartheta_0x_1,\dots,\vartheta_0x_n,y_1+te_0,
\dots,y_{n'}+te_0$ of the products of the covariances of the pairs, and
$\mathbb E[\vartheta_0(F)]\,\mathbb E[\tau_tF']$ is the sum over those pairings
that pair no $\vartheta_0x_i$ with a $y_j+te_0$. Every pairing in the
difference contains a cross pair, whose covariance $C(y_j+te_0-\vartheta_0x_i)$
has time component $y_{j,0}+t+x_{i,0}-1\ge t$ in absolute value and is bounded
by $Ke^{-\epsilon t}$; the other factors are bounded by $C(0)$, and there
are at most $(n+n'-1)!!\le(n+n')!$ pairings. Hence
\[
\bigl|\mathbb E[\vartheta_0(F)\tau_tF']-\mathbb E[F]\,\mathbb E[F']\bigr|\le
(n+n')!\,K\,C(0)^{(n+n'-2)/2}\,e^{-\epsilon t},
\]
using the invariance of the law under $\vartheta_0$ and translations. By
linearity the same bound, with another constant, holds for arbitrary
polynomials $F=F'$ in finitely many $\xi_x$, $x\in\Lambda^+_{0,1}$, which is
the condition (b) of Proposition~\ref{prop:gapdecay} with $m=\epsilon$; its
proof gives (a).
\end{proof}
% The anisotropic free field and the no-go theorem for rotation invariance.
\subsubsection{The anisotropic free field}

Fix $m>0$ and $\lambda>0$ with $\lambda\ne1$. For $\mu=0,\dots,3$ let
$c^{(\mu)}\in\R^4$ have $c^{(\mu)}_\mu=\lambda^2$ and $c^{(\mu)}_\nu=1$ for
$\nu\ne\mu$, and let $D_\mu$ be the diagonal matrix with entries
$(c^{(\mu)}_\nu)^{1/2}$, $\det D_\mu=\lambda$. Put $c_-=\min(1,\lambda^2)$,
$c_+=\max(1,\lambda^2)$ and
\[
m'=\frac m{\max(1,\lambda)} .
\]
Let $a_k\downarrow0$ and $k_0$ be such that $ma_k\le1$ and $a_k\le1$ for
$k\ge k_0$. For $k\ge k_0$ let $\phi^{(0)}_k,\dots,\phi^{(3)}_k$ be
independent centred Gaussian fields on $\Z^4$ with covariances
$C^{(\mu)}_k=C_{ma_k,c^{(\mu)}}$ of~\eqref{rN5:eq:green}, that is the Gibbs
fields of the nearest-neighbour actions $\frac12\sum_x[m^2a_k^2\phi_x^2+
\sum_\nu c^{(\mu)}_\nu(\phi_{x+e_\nu}-\phi_x)^2]$, and let
\[
\phi_k=\tfrac12\sum_{\mu=0}^3\phi^{(\mu)}_k,\qquad C_k=\tfrac14\sum_\mu
C^{(\mu)}_k ,
\]
so that $\phi_k$ is a centred Gaussian field with covariance $C_k(x-y)$.
The tuple $(\phi^{(\mu)}_k)_\mu$ is the Gibbs field of the sum over $\mu$ of
these actions, a nearest-neighbour action for a field with values in
$\R^4$; a permutation of the coordinate axes permutes its components. The field
$\phi_k$ itself is a linear image of it, not a Gibbs field of a
nearest-neighbour action (its precision $C_k^{-1}$ is not of finite range).
Let $\omega_k$ be the expectation, $\Phi_x=\phi_k(x)$ (so $r=0$ and
$c_k=\omega_k(\Phi_0)=0$), $Z_k=a_k^{-1}$ and $A_k(f)=\sum_xa_k^4f(a_kx)
\phi_k(x)$.

On the continuum side let, for $y\in\R^4\setminus\{0\}$,
\begin{equation}\label{rN5:eq:Gm}
\begin{gathered}
G_m(y)=\int_0^\infty\frac{e^{-|y|^2/(4s)-m^2s}}{(4\pi s)^2}\,ds,\\
W^{(\mu)}(y)=\lambda^{-1}G_m(D_\mu^{-1}y),\qquad W=\tfrac14\sum_\mu W^{(\mu)} .
\end{gathered}
\end{equation}

\begin{lemma}\label{rN5:cont}
\textup{(i)} $G_m$ is positive and real analytic on $\R^4\setminus\{0\}$,
\[
G_m(y)\le\frac1{4\pi^2|y|^2},\qquad G_m(y)\le\frac{(5+m|y|)\,e^{-m|y|}}
{\pi^2|y|^2},
\]
$G_m\in L^1(\R^4)$, $\int G_m=m^{-2}$, and $\widehat{G_m}(q)=\int e^{-iq\cdot
y}G_m(y)\,dy=(m^2+|q|^2)^{-1}$.

\textup{(ii)} Let $E(\vec q)=(m^2+|\vec q|^2)^{1/2}$. For $y_0\ne0$,
\[
G_m(y)=\int_{\R^3}e^{i\vec q\cdot\vec y}\,\frac{e^{-E(\vec q)|y_0|}}
{2E(\vec q)}\,\frac{d^3q}{(2\pi)^3},
\]
and $W^{(\mu)}=W_{c^{(\mu)}}$ in the notation of
Lemma~\textup{\ref{rN5:green}(iv)}; the same holds with any coordinate
direction in the role of $0$.

\textup{(iii)} $\widehat{W^{(\mu)}}(q)=(m^2+\sum_\nu c^{(\mu)}_\nu
q_\nu^2)^{-1}$ and $0<\widehat W\le m^{-2}$. With $K_W=5\lambda^{-1}c_+
\max(1,\lambda^{-1})\max(1,m)/\pi^2$,
\[
0<W(y)\le K_W\,\frac{1+|y|}{|y|^2}\,e^{-m'|y|} .
\]
$W$ is invariant under $B_4$ and real analytic on $\R^4\setminus\{0\}$.
\end{lemma}

\begin{proof}
(i) Positivity is clear. The integral defining $G_m$ converges uniformly on
compact subsets of $\{w\in\C:\operatorname{Re}w>0\}$ when $|y|^2$ is replaced
by $w$, since $|e^{-w/(4s)}|=e^{-\operatorname{Re}w/(4s)}$; so $G_m$ is a
holomorphic function of $|y|^2$, which gives real analyticity. Dropping
$e^{-m^2s}$ and using $\int_0^\infty s^{-2}e^{-a/s}ds=a^{-1}$ gives the
first bound. For the second write
\[
\frac{|y|^2}{4s}+m^2s=m|y|+\frac{(|y|-2ms)^2}{4s}.
\]
On $s\le|y|/(4m)$ we have $(|y|-2ms)^2\ge|y|^2/4$, and
\[
\int_0^\infty\frac{e^{-|y|^2/(16s)}}{(4\pi s)^2}\,ds=\frac1{\pi^2|y|^2}.
\]
On $s\ge|y|/(4m)$ we have
$(4\pi s)^{-2}\le m^2/(\pi^2|y|^2)$; the factor $e^{-(|y|-2ms)^2/(4s)}$ is
at most $1$ on $[|y|/(4m),|y|/m]$, an interval of length at most $|y|/m$,
and at most $e^{-m^2s/4}$ for $s\ge|y|/m$, where $2ms-|y|\ge ms$. So this
part is at most $(m^2/\pi^2|y|^2)(|y|/m+4/m^2)$, and the bound follows. By
Fubini, $\int G_m=\int_0^\infty e^{-m^2s}ds=m^{-2}$, and the Gaussian
integral $\int e^{-iq\cdot y}(4\pi s)^{-2}e^{-|y|^2/(4s)}dy=e^{-s|q|^2}$
gives $\widehat{G_m}(q)=\int_0^\infty e^{-(m^2+|q|^2)s}ds$.

(ii) Factor the heat kernel in the time and space variables,
\[
\frac{e^{-|y|^2/(4s)}}{(4\pi s)^2}=\frac{e^{-y_0^2/(4s)}}{(4\pi s)^{1/2}}
\int_{\R^3}e^{i\vec q\cdot\vec y-s|\vec q|^2}\frac{d^3q}{(2\pi)^3}.
\]
By Fubini (the integrands are integrable in absolute value, by the
evaluation that follows),
\[
G_m(y)=\int_{\R^3}e^{i\vec q\cdot\vec y}\int_0^\infty\frac{e^{-y_0^2/(4s)
-E^2s}}{(4\pi s)^{1/2}}\,ds\,\frac{d^3q}{(2\pi)^3}.
\]
With $s=u^2/E^2$,
\[
\int_0^\infty\frac{e^{-a^2/(4s)-E^2s}}{s^{1/2}}\,ds=\frac2EJ\Bigl(\frac{aE}
2\Bigr),
\]
where $J(b)=\int_0^\infty e^{-u^2-b^2/u^2}du$. The substitution $u\mapsto b/u$
in
\[
J'(b)=-\int_0^\infty\frac{2b}{u^2}\,e^{-u^2-b^2/u^2}du
\]
gives $J'=-2J$, so $J(b)=J(0)e^{-2b}=\frac{\sqrt\pi}2e^{-2b}$ and the inner
integral is $e^{-E|y_0|}/(2E)$. For $W^{(\mu)}$ put $c=c^{(\mu)}$
and substitute $y'=D_\mu^{-1}y$, $q'_i=c_i^{1/2}q_i$: then $d^3q'=(\lambda
/\sqrt{c_0})\,d^3q$, $\vec q'\cdot\vec y'=\vec q\cdot\vec y$, $E(\vec q')=
\sqrt{c_0}\,\Omega(\vec q)$ and $|y'_0|=|y_0|/\sqrt{c_0}$, which gives
\[
W^{(\mu)}(y)=\int_{\R^3}e^{i\vec q\cdot\vec y}\,\frac{e^{-\Omega(\vec q)
|y_0|}}{2c_0\Omega(\vec q)}\,\frac{d^3q}{(2\pi)^3}=W_c(y).
\]
Since $G_m$ is rotation invariant, the same holds for every coordinate
direction.

(iii) By $y=D_\mu y'$, $\widehat{W^{(\mu)}}(q)=\lambda^{-1}\int e^{-iq\cdot
y}G_m(D_\mu^{-1}y)\,dy=\widehat{G_m}(D_\mu q)$, and $|D_\mu q|^2=\sum_\nu
c^{(\mu)}_\nu q_\nu^2$. The bounds follow from (i), since $|D_\mu^{-1}y|$ lies
between $|y|/\max(1,\lambda)$ and $|y|\max(1,\lambda^{-1})$ and $m|D_\mu^{-1}y|
\ge m'|y|$.
A permutation of the coordinates permutes the $W^{(\mu)}$ and a sign change
fixes each of them, so $W$ is $B_4$-invariant.
\end{proof}

The cube condition holds for these fields on tori, by the argument of
Theorem~\ref{thm:ffscalar}. For $L\ge n+1$, $n\ge2$ and $M>0$ let
$\mathbb T_L=(\Z/L\Z)^4$, let $\Delta_\nu$ be the lattice Laplacian of
$\mathbb T_L$ in direction $\nu$, $Q^{(\mu)}=M^2-\sum_\nu c^{(\mu)}_\nu
\Delta_\nu$, let $(\phi^{(\mu)})_\mu$ be the centred Gaussian field on
$\mathbb T_L$ with precision $\bigoplus_\mu Q^{(\mu)}$ (independent
components), $\phi=\frac12\sum_\mu\phi^{(\mu)}$, and for $x\in\mathbb T_L$
let $V_x=x+\{0,\dots,n-1\}^4$ and $\mathcal P_x$ the conditional expectation
given $(\phi_y)_{y\notin V_x}$.

\begin{lemma}[the cube condition]\label{rN5:bp}
For every real $f\in L^2$ measurable with respect to $\phi$,
\[
\varpi\Var(f)\le\frac1{n^4}\sum_{x\in\mathbb T_L}\E\bigl(f-\mathcal P_xf
\bigr)^2,
\]
where $s=nM$ and
\[
\varpi=\Bigl(\frac56\Bigr)^4\frac{s^2}{s^2+10(3+\lambda^2)} .
\]
In particular, if $\ell>0$, $M=ma_k$ and $n=n_k=\lceil\ell/a_k\rceil\ge2$,
then $s\ge m\ell$ and $\mathrm{BP}(n_k,\varpi_\ell)$ holds on every such
torus with $\varpi_\ell=(5/6)^4(m\ell)^2/((m\ell)^2+10(3+\lambda^2))$,
independent of $k$ and $L$.
\end{lemma}

\begin{proof}
Let first $f$ be a function of the tuple, $N=4L^4$, $Q=\bigoplus_\mu
Q^{(\mu)}$, and let $\mathcal P'_x$ be the conditional expectation given the
tuple outside $V_x$, that is, given the coordinates outside $V'_x=V_x\times
\{0,\dots,3\}$. With the parabolic profile $\tau$ of
Lemma~\ref{lem:ffparabola}, $\chi_x$ as before Lemma~\ref{lem:ffpartition}
and $v=(v^{(\mu)})_\mu\in\R^N$, put $v_x=(\chi_xv^{(\mu)})_\mu$, supported in
$V'_x$, with $\sum_xv_x=v$. The proof of Lemma~\ref{lem:ffpartition} gives,
for each direction $\nu$ and each $w\in\R^{\mathbb T_L}$,
$\sum_x\|\nabla_\nu(\chi_xw)\|^2=S_2^3[S_{11}\|\nabla_\nu w\|^2+D_2\|w\|^2]$.
Multiplying by $c^{(\mu)}_\nu$, summing over $\nu$, adding $M^2\sum_x\|\chi_x
w\|^2=M^2S_2^4\|w\|^2$ and using $S_{11}\le S_2$ and $\sum_\nu c^{(\mu)}_\nu=
3+\lambda^2$,
\[
\sum_x\ip{\chi_xw}{Q^{(\mu)}\chi_xw}\le C_0\ip w{Q^{(\mu)}w},
\]
where
\[
C_0=S_2^3\Bigl(S_2+\frac{(3+\lambda^2)D_2}{M^2}\Bigr),
\]
for every $\mu$, hence $\sum_x\ip{v_x}{Qv_x}\le C_0\ip v{Qv}$.
Lemma~\ref{lem:fflions} with $T=\R^N$ and the coordinate sets $V'_x$ gives
$\varpi\Var(f)\le n^{-4}\sum_x\E(f-\mathcal P'_xf)^2$ with $\varpi=1/(C_0
n^4)$, and $C_0n^4=\alpha_n^4+(3+\lambda^2)\alpha_n^3\beta_ns^{-2}=\alpha_n^4(1+
10\rho_n(3+\lambda^2)s^{-2})$; by $\alpha_n\le\frac65$ and $\rho_n<1$
(Lemma~\ref{lem:ffparabola}), $1/(C_0n^4)$ is at least the stated $\varpi$.
Now let $f$ be a function of $\phi$. Since $\phi_y$ is a function of the
tuple at $y$, the $\sigma$-algebra of $(\phi_y)_{y\notin V_x}$ is contained in
that of the tuple outside $V_x$, so $\mathcal P_xf$ is measurable for the
latter and $\E(f-\mathcal P'_xf)^2\le\E(f-\mathcal P_xf)^2$, $\mathcal P'_x$
being the orthogonal projection onto the functions measurable for the larger
$\sigma$-algebra. This gives the claim. For $n=\lceil\ell/a_k\rceil$,
$s=nma_k\ge m\ell$, and $s^2/(s^2+c)$ increases with $s$.
\end{proof}

As $L\to\infty$ the covariance of $\phi^{(\mu)}$ on $\mathbb T_L$ converges
to $C_{M,c^{(\mu)}}$, being the Riemann sum over $(2\pi/L)\Z^4\cap T^4$ of
the continuous integrand of~\eqref{rN5:eq:green}; so the states $\omega_k$
below are the limits of these torus states with $M=ma_k$.

\subsubsection{The no-go theorem}

\begin{theorem}[the anisotropic free field]\label{rN5:nogo}
Let $m$, $\lambda$, $a_k$, $\phi_k$, $\omega_k$, $Z_k$, $A_k$ and $W$ be as
above, and $k\ge k_0$.

\textup{(i)} \emph{Lattice.} Each $\omega_k$ is invariant under the
translations of $\Z^4$ and under $B_4$, and it is reflection positive for
the site reflection $\theta_\mu$ and the link reflection $\vartheta_\mu$ in
every direction $\mu$. The transfer matrix of $\omega_k$ in every direction
has spectrum on $\Omega_k^\perp$ in $[0,e^{-\epsilon_k}]$ with
\[
\epsilon_k\ge a_km'\Bigl(1+\frac{m^2a_k^2}{4c_-}\Bigr)^{-1/2} .
\]
Consequently, for every $m_*<m'$ there is $k_1$ such that
Lemma~\textup{\ref{lem:limitElattice}(iii)} holds for $k\ge k_1$ with this
$m_*$, and so do the conclusions of Theorems~\textup{\ref{thm:limitEsep}},
\textup{\ref{thm:limitEtempered}} and Corollary~\textup{\ref{cor:limitEexact}}
\textup{(}Remark~\textup{\ref{rem:limitEscope}}\textup{)}. The torus states
whose limit is $\omega_k$ satisfy $\mathrm{BP}(n_k,\varpi_\ell)$ for every
$\ell>0$, with $n_k=\lceil\ell/a_k\rceil$ and $\varpi_\ell$ independent of
$k$ \textup{(}Lemma~\textup{\ref{rN5:bp}}\textup{)}.

\textup{(ii)} \emph{Moments.} For every $n$ and all $f_1,\dots,f_n\in C_c
(\R^4)$, with no separation of the supports,
\[
\begin{split}
\lim_{k\to\infty}\omega_k\Bigl(\prod_{j=1}^nZ_kA_k(f_j)\Bigr)
&=S_n(f_1,\dots,f_n)\\
&=\sum_{\text{pairings }\pi}\ \prod_{\{i,j\}\in\pi}S_2(f_i,f_j),
\end{split}
\]
where $S_2(f,g)=\iint f(x)g(y)W(x-y)\,dx\,dy$ and the sum runs over the pairings of $\{1,\dots,n\}$ \textup{(}zero for odd
$n$\textup{)}. For all $k\ge k_0$ and $f,g\in C_c(\R^4)$,
\begin{equation}\label{rN5:eq:l2}
Z_k^2\bigl|\omega_k(A_k(f)A_k(g))\bigr|\le m^{-2}\|f\|_{2,k}\|g\|_{2,k},
\end{equation}
where $\|f\|_{2,k}^2=a_k^4\sum_x|f(a_kx)|^2$.
Hence \textup{(M)}, \textup{(M$_{\mathrm s}$)}, \textup{(G$_{\mathrm s}$)}
with $K=4/m$, $\beta=\frac12$, $q=2$, $p=0$, and \textup{(U$_{\mathrm s}$)}
hold, and each $S_n$ is a tempered distribution on $\R^{4n}$ with
\begin{equation}\label{rN5:eq:tempered}
|S_n(F)|\le\frac{(n-1)!!}{(4m)^n}\sup_{p\in\R^{4n}}(1+|p|^2)^{2n}|\hat F(p)|
\end{equation}
for $F\in\mathcal S(\R^{4n})$,
so that the Osterwalder--Schrader temperedness and linear growth
conditions~\cite{OS75} hold. Off the coincidence set
$\{x_i=x_j\text{ for some }i\ne j\}$ the $S_n$ are the real analytic
functions $\sum_\pi\prod_{\{i,j\}\in\pi}W(x_i-x_j)$, and $W>0$.

\textup{(iii)} \emph{Reflection positivity and gap.} For every $\mu$, $S(
\theta_\mu\bar P\cdot P)\ge0$ for every polynomial $P$ in the $A(f)$ with
$f\in C_c(\{x_\mu>0\})$. The Hamiltonian $H_\mu$ of the reconstruction of
Theorem~\textup{\ref{thm:limitEsep}(iv)} satisfies
\[
\sigma(H_\mu)\subset\{0\}\cup[m',\infty),\qquad\ker H_\mu=\C\Omega_\mu,
\qquad m'\in\sigma(H_\mu):
\]
the gap is exactly $m'$ in every direction, and the clustering of
Theorem~\textup{\ref{thm:limitEsep}(v)} holds with rate $m'$.

\textup{(iv)} \emph{Symmetry.} $G_S=B_4$. More precisely, $S_2(f\circ R^{-1},
g\circ R^{-1})=S_2(f,g)$ for all separated pairs forces $R\in B_4$. In
particular the limit displayed in \textup{(W$_{\mathrm s}$)} is nonzero for
some separated tuple, for every $R\in\mathrm O(4)\setminus B_4$, so
\textup{(W$_{\mathrm s}$)} fails, and the whole-sequence limit $S$ is not
invariant under any rotation outside $B_4$.
\end{theorem}

\begin{proof}
(i) The law of the tuple $(\phi^{(\mu)}_k)_\mu$ is invariant under
translations; a permutation $\sigma$ of the coordinates maps $\phi^{(\mu)}_k
\circ\sigma$ to a field with the law of $\phi^{(\sigma\mu)}_k$ (the
permutation acts on $c^{(\mu)}$ by permuting its entries), and a sign change
of a coordinate preserves each $C^{(\mu)}_k$; so the law of $\phi_k$ is
invariant under translations and $B_4$. Reflection positivity: by
Lemma~\ref{rN5:gaussrp} it suffices that $(x,y)\mapsto C_k(\theta_\mu x-y)$
is positive semidefinite on $\{x_\mu\ge0\}$ and $(x,y)\mapsto C_k(\vartheta_\mu
x-y)$ on $\{x_\mu\ge1\}$; as $C_k$ is a sum of the $C^{(\nu)}_k$, it suffices
to treat each $C^{(\nu)}_k$, and by the symmetry of~\eqref{rN5:eq:green} we
may take $\mu=0$. By~\eqref{rN5:eq:3d}, for $x_0,y_0\ge0$,
\[
C^{(\nu)}_k(\theta_0x-y)=\int_{T^3}\frac{\overline{e^{-\varepsilon x_0}
e^{-i\vec p\cdot\vec x}}\,e^{-\varepsilon y_0}e^{-i\vec p\cdot\vec y}}
{2c_0\sinh\varepsilon}\,\frac{d^3p}{(2\pi)^3},
\]
so that, for finitely supported $v$,
\[
\sum_{x,y}\bar v_xv_yC^{(\nu)}_k(\theta_0x-y)=\int_{T^3}\frac{\bigl|\sum_xv_x
e^{-\varepsilon x_0}e^{-i\vec p\cdot\vec x}\bigr|^2}{2c_0\sinh\varepsilon}\,
\frac{d^3p}{(2\pi)^3}\ge0;
\]
for $x_0,y_0\ge1$ the time component of $\vartheta_0x-y$ is
$1-x_0-y_0$, of absolute value $x_0+y_0-1$, and the same computation with an
extra factor $e^{\varepsilon}$ applies. By Lemma~\ref{rN5:green}(iii) applied
to each $C^{(\nu)}_k$ with $M=ma_k$, $|C_k(z)|\le K_3e^{-\epsilon_k|z_0|}$ with
\[
\begin{split}
\epsilon_k&=\min_\nu\operatorname{arccosh}\Bigl(1+\frac{m^2a_k^2}{2c^{(\nu)}_0}
\Bigr)\\
&\ge\frac{ma_k}{\sqrt{c_+}}\Bigl(1+\frac{m^2a_k^2}{4c_-}\Bigr)^{-1/2},
\end{split}
\]
and $m/\sqrt{c_+}=m'$;
Lemma~\ref{rN5:gaussgap} gives the spectral bound in direction $0$, and the
other directions follow from the $B_4$-invariance. Given $m_*<m'$, choose $k_1
\ge k_0$ with $\epsilon_k\ge m_*a_k$ for $k\ge k_1$; then the proof of
Lemma~\ref{lem:limitElattice}(iii) applies verbatim, and by
Remark~\ref{rem:limitEscope} the hypothesis of Theorem~\ref{thm:avgtransfer}
is used in Section~\ref{ssec:limitE} only through that lemma.

(ii) Let first $f,g\in C^\infty_c(\R^4)$ be real, $F_k(x)=f(a_kx)$, $G_k(x)=
g(a_kx)$ and $\hat F_k(p)=\sum_xF_k(x)e^{-ip\cdot x}$. By Parseval on $T^4$,
\[
\begin{split}
Z_k^2\omega_k(A_k(f)A_k(g))&=a_k^6\sum_{x,y}F_k(x)G_k(y)C_k(x-y)\\
&=a_k^6\int_{T^4}\overline{\hat F_k(p)}\hat G_k(p)\hat C_k(p)\,
\frac{d^4p}{(2\pi)^4},
\end{split}
\]
where $0<\hat C_k(p)=\frac14\sum_\mu(m^2a_k^2+\sum_\nu c^{(\mu)}_\nu2(1-\cos
p_\nu))^{-1}\le(ma_k)^{-2}$. The Cauchy--Schwarz inequality and Parseval
give~\eqref{rN5:eq:l2}, which extends to $f,g\in C_c(\R^4)$ by uniform
approximation (both sides are continuous in $\|\cdot\|_\infty$ on functions
supported in a fixed cube). Substituting $p=a_kq$,
\[
Z_k^2\omega_k(A_k(f)A_k(g))=\int_{|q_\nu|\le\pi/a_k}\overline{\hat f_k(q)}\,
\hat g_k(q)\,\hat W_k(q)\,\frac{d^4q}{(2\pi)^4},
\]
where
\[
\hat W_k(q)=\frac14\sum_\mu\Bigl(m^2+\sum_\nu c^{(\mu)}_\nu\frac{2(1-\cos
a_kq_\nu)}{a_k^2}\Bigr)^{-1},
\]
and $\hat f_k(q)=a_k^4\hat F_k(a_kq)=\sum_xa_k^4f(a_kx)e^{-iq\cdot a_kx}$, a
Riemann sum converging to $\hat f(q)$ for each $q$. For $N\in\mathbb N$ and
each $\nu$, $(1-e^{-ia_kq_\nu})^N\hat f_k(q)=\sum_xa_k^4(\nabla_\nu^Nf)(a_kx)
e^{-iq\cdot a_kx}$ with $\nabla_\nu f(y)=f(y)-f(y-a_ke_\nu)$, and $|\nabla_\nu^N
f|\le a_k^N\|\partial_\nu^Nf\|_\infty$; since $|1-e^{-i\theta}|\ge2|\theta|/
\pi$ on $[-\pi,\pi]$, $|\hat f_k(q)|\le(\pi/2)^N\|\partial_\nu^Nf\|_\infty
|\operatorname{supp}f+[-6,6]^4|\,|q_\nu|^{-N}$ for $|a_kq_\nu|\le\pi$ and
$N\le5$ \textup{(}the support of $\nabla^N_\nu f$ lies within $\ell^\infty$
distance $5a_k\le5$ of $\operatorname{supp}f$\textup{)}, which
with $N=0$ and $N=5$ gives $|\hat f_k(q)|\le C_f(1+|q|_\infty)^{-5}$ uniformly
in $k$. The factor $\hat W_k(q)$ is at most $m^{-2}$ and converges to
$\hat W(q)$. By dominated convergence and Lemma~\ref{rN5:cont}(iii),
\[
\begin{split}
\lim_kZ_k^2\omega_k(A_k(f)A_k(g))&=\int\overline{\hat f}\hat g\hat W\,\frac
{d^4q}{(2\pi)^4}\\
&=\iint f(x)g(y)W(x-y)\,dx\,dy ,
\end{split}
\]
the last step by Fubini and $\widehat{W}=\hat W$ with $W\in L^1$. For
$f,g\in C_c(\R^4)$ the convergence follows from~\eqref{rN5:eq:l2}, the bound
$|S_2(f,g)|\le m^{-2}\|f\|_2\|g\|_2$ (from $\hat W\le m^{-2}$) and uniform
approximation by smooth functions supported in a fixed cube. Since $\phi_k$ is
Gaussian, $\omega_k(\prod_jZ_kA_k(f_j))$ is the sum over pairings of the
products of $Z_k^2\omega_k(A_k(f_i)A_k(f_j))$ (Isserlis~\cite{Iss18}), which
gives the convergence of all moments and the Gaussian form of the limit.
For $f_j$ supported in a cube $Q$ of side $L$, $\|f_j\|_{2,k}^2\le(L+2a_k)^4
\|f_j\|_\infty^2\le16(1+L)^4\|f_j\|_\infty^2$, and there are $(n-1)!!\le
(n!)^{1/2}$ pairings, so $|\omega_k(\prod_jZ_kA_k(f_j))|\le(n!)^{1/2}(4(1+L)^2/
m)^n\prod_j\|f_j\|_\infty$ for all $k\ge k_0$: this is (M), (M$_{\mathrm s}$)
and (G$_{\mathrm s}$) with the stated constants, and (U$_{\mathrm s}$) holds
because the whole sequence converges. For~\eqref{rN5:eq:tempered} let $F\in
\mathcal S(\R^{4n})$, $n$ even, and $\pi$ a pairing. In the variables
$u_{ij}=x_i-x_j$ and $x_j$, one pair $\{i,j\}$, $i<j$, of $\pi$ at a time,
$\int F(x)\prod_{\{i,j\}\in\pi}W(x_i-x_j)\,dx=\int_{(\R^4)^{n/2}}\prod W(u_{ij})
\,g(u)\,du$ with $g$ the integral of $F$ over the $x_j$, a Schwartz function
with $\hat g(q)=\hat F(\iota_\pi q)$, where $\iota_\pi q$ puts $q_{ij}$ at
the place $i$ and $-q_{ij}$ at the place $j$. As $\prod W(u_{ij})\in L^1$ and
$\hat g\in L^1$, Fubini and Fourier inversion give $\int\prod W(u_{ij})g(u)
\,du=(2\pi)^{-2n}\int\prod\hat W(-q_{ij})\hat F(\iota_\pi q)\,dq$, so this is at
most $(2\pi)^{-2n}m^{-n}\int_{\R^{2n}}|\hat F(\iota_\pi q)|\,dq$; $|\iota_\pi q|^2=2|q|^2$
and $(1+2|q|^2)^{2n}\ge\prod_{i=1}^{2n}(1+q_i^2)$ give $\int_{\R^{2n}}(1+2|q|^2)
^{-2n}dq\le\pi^{2n}$, and~\eqref{rN5:eq:tempered} follows. For the linear
growth condition note that $(1+|p|^2)^{2n}\hat F$ is the Fourier transform of
$(1-\Delta)^{2n}F$, so the supremum in~\eqref{rN5:eq:tempered} is at most
$\|(1-\Delta)^{2n}F\|_{L^1(\R^{4n})}$; expanding $(1-\Delta)^{2n}=\sum_j\binom
{2n}j(-\sum_{i=1}^{4n}\partial_i^2)^j$ gives $|(1-\Delta)^{2n}F(x)|\le(1+4n)^{2n}
\sup_{|\gamma|\le4n}|\partial^\gamma F(x)|$, and $\int_{\R^{4n}}(1+|x|^2)^{-(4n+1)
/2}dx\le\bigl(\int_{\R^2}(1+|w|^2)^{-1-1/(4n)}d^2w\bigr)^{2n}=(4\pi n)^{2n}$,
since $(1+|x|^2)^{(4n+1)/2}\ge\prod_{i=1}^{2n}(1+|x_{(i)}|^2)^{(4n+1)/(4n)}$
for the $2n$ pairs $x_{(i)}\in\R^2$ of coordinates. With the norms $|F|_w=
\sup_{x,\,|\gamma|\le w}(1+|x|^2)^{w/2}|\partial^\gamma F(x)|$ of
Theorem~\ref{thm:limitEtempered} this gives $|S_n(F)|\le\sigma_n|F|_{5n}$, $\sigma_n=
(n!)^{1/2}(4m)^{-n}(5n)^{2n}(4\pi n)^{2n}\le(100\pi^2e^4/m)^n(n!)^{9/2}$
\textup{(}by $n^n\le e^nn!$\textup{)},
the linear growth condition of~\cite{OS75} with $s=5$. Real analyticity
off the coincidence set is Lemma~\ref{rN5:cont}(iii).

(iii) Let $f,g\in C_c(\{x_\mu>0\})$. By Lemma~\ref{rN5:cont}(ii) with $\mu$ as
time direction, for $x_\mu,y_\mu>0$,
\[
W^{(\nu)}(\theta_\mu x-y)=\int_{\R^3}\frac{\overline{e^{-\Omega x_\mu}
e^{-i\vec q\cdot\vec x}}\,e^{-\Omega y_\mu}e^{-i\vec q\cdot\vec y}}
{2c^{(\nu)}_\mu\Omega}\,\frac{d^3q}{(2\pi)^3},
\]
so $(f,g)\mapsto\iint f(x)g(y)
W(\theta_\mu x-y)$ is a positive semidefinite kernel on $C_c(\{x_\mu>0\})$.
The Gaussian family $S$ consists of the moments of the centred Gaussian
process $(\Phi(f))_{f\in C_c(\R^4)}$ with covariance $S_2$, which exists by
Kolmogorov's theorem since $S_2$ is positive semidefinite
($\hat W>0$), and whose law is invariant under $\theta_\mu$; the proof of
Lemma~\ref{rN5:gaussrp}, with sites replaced by test functions, gives $S(\theta_\mu\bar P\cdot P)\ge0$. Clustering: for $f,g
\in C_c(\{x_\mu\ge\delta\})$ and $t\ge0$,
\[
S_2(\theta_\mu\bar f,\tau_{te_\mu}g)=\iint\bar f(x)g(y)W(\theta_\mu x-y-te_\mu)
\,dx\,dy,
\]
where the $\mu$-component of the
argument has absolute value $x_\mu+y_\mu+t\ge2\delta+t$. For $|z_\mu|\ge
2\delta$ Lemma~\ref{rN5:cont}(iii) gives $|W(z)|\le K_W(1+|z|)|z|^{-2}e^{-m'
|z|}\le K_\delta e^{-m'|z_\mu|}$ with $K_\delta=K_W(2\delta)^{-2}(1+2\delta)$,
since $|z|\ge|z_\mu|\ge2\delta$. So each cross covariance in the Wick
expansion of $S(\theta_\mu\bar P\cdot\tau_{te_\mu}P')-\overline{S(P)}S(P')$ is
bounded by $K_\delta\|f\|_1\|g\|_1e^{-m't}$, and as in the proof of
Lemma~\ref{rN5:gaussgap}, $|S(\theta_\mu\bar P\cdot\tau_{te_\mu}P)-|S(P)|^2|
\le C_Pe^{-m't}$ for every polynomial $P$ in the $A(f)$ with $f\in C_c(\{x_\mu
>0\})$, in particular for every separated polynomial at positive $\mu$-times.
In the reconstruction of Theorem~\ref{thm:limitEsep}(iii), (iv) these numbers
are
\[
\ip{\Psi'_P}{e^{-tH_\mu}\Psi'_P}=\int e^{-tE}d\nu_P(E),
\]
with $\Psi'_P=\Psi_\mu(P)-\ip{\Omega_\mu}{\Psi_\mu(P)}\Omega_\mu$ and $\nu_P$ its spectral
measure; for $E_0<m'$, $\nu_P([0,E_0])e^{-tE_0}\le C_Pe^{-m't}$ for all $t$
forces $\nu_P([0,E_0])=0$, and the $\Psi'_P$, $P$ separated, are dense in
$\Omega_\mu^\perp$ (the $\Psi_\mu(P)$ are dense in $\mathcal H_\mu$ by
construction), so $1_{[0,E_0]}(H_\mu)$ vanishes on $\Omega_\mu^\perp$:
$\sigma(H_\mu)\subset\{0\}\cup[m',\infty)$ and $\ker H_\mu=\C\Omega_\mu$.
For the lower bound let $\nu$ be
an index with $c^{(\nu)}_\mu=c_+$, let $f(x)=g(x_\mu)h(\vec x)$ with $0\le g\in
C_c((0,\infty))$, $g\ne0$, $h\in C_c(\R^3)$ real with $\hat h(0)\ne0$, and
$\Psi=\Psi_\mu(A(f))\perp\Omega_\mu$. By Lemma~\ref{rN5:cont}(ii),
\[
\begin{split}
\ip\Psi{e^{-tH_\mu}\Psi}&=S_2(\theta_\mu f,\tau_{te_\mu}f)\\
&=\frac14\sum_{\nu'}\int_{\R^3}\frac{|\hat h(\vec q)|^2\,
\tilde g(\Omega_{\nu'})^2}{2c^{(\nu')}_\mu\Omega_{\nu'}}\,
e^{-t\Omega_{\nu'}}\,\frac{d^3q}{(2\pi)^3},
\end{split}
\]
where $\tilde g(\Omega)=\int_0^\infty g(s)e^{-\Omega s}ds>0$ and
$\Omega_{\nu'}=\Omega^{(\nu')}(\vec q)$ with $\Omega^{(\nu')}(\vec q)=((m^2+\sum_{i\ne\mu}c^{(\nu')}_iq_i^2)/c^{(\nu')}_\mu
)^{1/2}\ge m/\sqrt{c_+}=m'$, with equality for $\nu'=\nu$ at $\vec q=0$. So
$\ip\Psi{e^{-tH_\mu}\Psi}=\int e^{-t\Omega}d\rho(\Omega)$ for a positive
measure $\rho$ with $\rho([m',m'+\eta])>0$ for every $\eta>0$. If $\sigma(H_\mu)
\cap[m',m'+\eta)=\emptyset$, the spectral measure of $\Psi$ gives $\ip\Psi
{e^{-tH_\mu}\Psi}\le\|\Psi\|^2e^{-(m'+\eta)t}$, while the right side is at
least $\rho([m',m'+\eta/2])e^{-(m'+\eta/2)t}$, a contradiction for large
$t$. Hence $m'\in\sigma(H_\mu)$. The clustering statement of
Theorem~\ref{thm:limitEsep}(v) with rate $m'$ follows from its proof, which
uses only the spectral bound $\sigma(H_\mu)\subset\{0\}\cup[m_*,\infty)$, now
available with $m_*=m'$.

(iv) Suppose $S_2(f\circ R^{-1},g\circ R^{-1})=S_2(f,g)$ for all separated
pairs. Then $\iint f(x)g(y)[W(R(x-y))-W(x-y)]\,dx\,dy=0$ whenever the supports
of $f,g$ are disjoint, and since $W$ is continuous off $0$, $W\circ R=W$ on
$\R^4\setminus\{0\}$. As $W\in L^1(\R^4)$, $\hat W$ is continuous and $\hat W
(Rq)=\hat W(q)$ for all $q\in\R^4$. Fix a unit vector $v\in\R^4$ and put
$\alpha_\mu(v)=1+(\lambda^2-1)v_\mu^2\ge c_->0$. By Lemma~\ref{rN5:cont}(iii),
for $s\in\R$,
\[
\hat W(sv)=\frac14\sum_{\mu=0}^3\frac1{m^2+s^2\alpha_\mu(v)}=:\Phi_v(s^2),
\]
and $\Phi_v(u)=\frac14\sum_{\alpha}n_\alpha(m^2+u\alpha)^{-1}$, the sum over
the distinct values $\alpha$ of $\alpha_0(v),\dots,\alpha_3(v)$ with their
multiplicities $n_\alpha\ge1$, is a rational function of $u$ whose poles are
exactly the points $u=-m^2/\alpha$ (the residue there is $n_\alpha/(4\alpha)
\ne0$, and distinct $\alpha$ give distinct poles). Since $\Phi_{Rv}=\Phi_v$ on
$[0,\infty)$, these rational functions coincide, so they have the same poles:
$\{\alpha_\mu(Rv)\}_\mu=\{\alpha_\mu(v)\}_\mu$ as sets, and as $t\mapsto1+(
\lambda^2-1)t$ is injective, $\{(Rv)_\mu^2\}_\mu=\{v_\mu^2\}_\mu$ as sets. For
$v=e_\nu$ this set is $\{0,1\}$, so every coordinate of $Re_\nu$ is $0$ or
$\pm1$, and as $|Re_\nu|=1$ exactly one coordinate is $\pm1$: $R$ is a
signed permutation matrix, $R\in B_4$. Conversely $B_4\subset G_S$ by
Lemma~\ref{rN5:cont}(iii) and (ii). Since the $S_n$ are Wick sums of $S_2$,
$G_S$ is the invariance group of $S_2$, which is $B_4$. Because the limit
exists along the whole sequence, the limit in (W$_{\mathrm s}$) vanishes for
$R$ and every separated tuple if and only if $R\in G_S$; $B_4$ is finite,
so it contains no element of infinite order.
\end{proof}
% Consequences for rotation invariance, and nontriviality.
\begin{remark}[what a rotation theorem must assume]\label{rN5:rem:nogo}
The states of Theorem~\ref{rN5:nogo} are linear images of Gibbs states of a
nearest-neighbour action, translation and $B_4$ invariant, reflection
positive for both reflection types in all four directions, with transfer
matrices whose physical gap is bounded below uniformly in $k$, and they
satisfy the cube condition on tori uniformly at fixed physical cube size;
the whole sequence of moments converges, with (M), (M$_{\mathrm s}$),
(G$_{\mathrm s}$) and (U$_{\mathrm s}$); the limit is reflection positive in
all four directions, has the same gap $m'$ in all of them, clusters
exponentially, satisfies the Osterwalder--Schrader temperedness and linear
growth conditions, has positive, real analytic Schwinger functions off the
coincidence set and a two-point kernel with the free-field short-distance
law $|y|^{-2}$; and yet $G_S=B_4$. So every hypothesis of
Section~\ref{ssec:limitE} other than (W$_{\mathrm s}$) itself, read for
fields at the sites, is compatible with the failure of (W$_{\mathrm s}$),
and Corollary~\ref{cor:limitEexact} cannot be improved by dropping
(W$_{\mathrm s}$) from its list, nor by replacing it with any of the
properties above. The restoration of rotation invariance must therefore be
derived from a property of the Wilson measure that these examples lack,
such as the Symanzik-type conditions (SB$_n$), (RI$_n$) of
Theorem~\ref{thm:limitErot}. Four components are used because $B_4$
permutes them: for the field of a single component with a $B_4$-invariant
nearest-neighbour action $\frac12\sum_x[M^2\xi_x^2+c\sum_\nu(\xi_{x+e_\nu}-
\xi_x)^2]$ the substitution of the proof of Lemma~\ref{rN5:cont}(ii) gives
the limit kernel $W_{(c,c,c,c)}(y)=c^{-2}G_m(y/\sqrt c)$, which is radial. The isotropic case
$\lambda=1$ of the construction is the standard massive free lattice field,
for which the cube condition holds uniformly at fixed physical cube size
(Theorem~\ref{thm:ffscalar}) and $G_S=\mathrm O(4)$, since $G_m$ is radial:
the cube condition, the gap, $\mathrm O(4)$ invariance and
(U$_{\mathrm s}$) are all compatible with a Gaussian limit, so
nontriviality is a separate input as well.
\end{remark}

\begin{corollary}[infinitely many rotations suffice]\label{rN5:infmany}
Under the hypothesis of Theorem~\textup{\ref{thm:avgtransfer}} and
\textup{(M$_{\mathrm s}$)}, let $S$ be a limit of
Theorem~\textup{\ref{thm:limitEsep}}. If $S$ is invariant under infinitely
many elements of $\mathrm O(4)$, then $G_S=\mathrm O(4)$. In particular the
element of infinite order in \textup{(W$_{\mathrm s}$)} may be replaced by
infinitely many distinct rotations of finite order, for instance the
rotations by $2\pi/N$ in one fixed plane for infinitely many $N$: if the
displayed limit in \textup{(W$_{\mathrm s}$)} is zero for each of them, then
every limit $S$ along a subsequence has $G_S=\mathrm O(4)$.
\end{corollary}

\begin{proof}
$G_S$ is a closed subgroup containing $B_4$ which is finite or $\mathrm O(4)$
(Theorem~\ref{thm:limitEsep}(vii)); an infinite one is $\mathrm O(4)$. If the
limit in (W$_{\mathrm s}$) vanishes for $R$, every subsequential limit $S$ is
$R$-invariant on separated tuples, so $R\in G_S$.
\end{proof}

\subsubsection{Triviality: what the gap excludes and what it does not}

For a gauge-invariant composite field such as the plaquette sum $\Phi$,
``the limit is Gaussian'' is not the right notion of triviality: Wick
polynomials of a free field are in general not Gaussian. For instance, for a
centred Gaussian field $\xi$ with pointwise positive covariance $G$, the
connected four-point function of ${:}\xi^2{:}$ is the sum over the $48$
connected Wick graphs on four vertices of degree two, each a four-cycle
contributing a product of four values of $G$, which is positive. The
appropriate notion is: the limit is the family of Schwinger functions of a
Wick polynomial of a free field. The gap of Theorem~\ref{thm:avgtransfer}
excludes the massless case.

\begin{proposition}[the gap excludes Wick polynomials of massless free
fields]\label{rN5:nomassless}
Assume the hypothesis of Theorem~\textup{\ref{thm:avgtransfer}} and
\textup{(M$_{\mathrm s}$)}, and let $S$ be a limit of
Theorem~\textup{\ref{thm:limitEsep}} along a subsequence.

\textup{(i)} Let $T$ be a continuous function on $\R^4\setminus\{0\}$ of the
form $T=\sum_{j=1}^Jh_j$ with $h_j$ homogeneous of degree $-d_j$, $0<d_1<
\dots<d_J$, and $C^1$ on $\R^4\setminus\{0\}$, not all $h_j$ zero. Then $S_2$
is not given by $T$: there is a separated pair $(f,g)$ with $S_2(f,g)\ne
\iint f(x)g(y)T(x-y)\,dx\,dy$.

\textup{(ii)} Let $\varphi$ be the massless free scalar field on $\R^4$,
with covariance $G_0(y)=(4\pi^2|y|^2)^{-1}$, or let $A_\mu$ be the free
Maxwell field, whose field strength $F_{\mu\nu}=\partial_\mu A_\nu-\partial_\nu
A_\mu$ has the covariance $\ip{F_{\mu\nu}(x)F_{\rho\sigma}(0)}=-(\delta_{\nu
\sigma}\partial_\mu\partial_\rho-\delta_{\nu\rho}\partial_\mu\partial_\sigma-
\delta_{\mu\sigma}\partial_\nu\partial_\rho+\delta_{\mu\rho}\partial_\nu
\partial_\sigma)G_0(x)$. Let $\psi$ be a Wick polynomial without constant term in
$\varphi$, or in the $F_{\mu\nu}$, and their derivatives, whose two-point
function does not vanish identically on $\R^4\setminus\{0\}$. Then $S_2$ is not the two-point
function of $\psi$ on separated pairs. This applies to $\varphi$, to
${:}\varphi^j{:}$ with $\ip{{:}\varphi^j{:}(x)\,{:}\varphi^j{:}(0)}=j!\,G_0(x)^j$,
to $F_{01}$ with
\[
\ip{F_{01}(x)F_{01}(0)}=\frac{4|x|^2-8(x_0^2+x_1^2)}{4\pi^2|x|^6},
\]
so that
\[
\ip{F_{01}(te_3)F_{01}(0)}=\frac1{\pi^2t^4},
\]
and to $\psi=\sum_{\mu<\nu}{:}F_{\mu\nu}F_{\mu\nu}{:}$, with $\ip{\psi(te_3)\psi(0)}
\ge2/(\pi^4t^8)$. The same holds for the Maxwell field of $\U(1)^{\dim G}$,
$\dim G$ independent copies.
\end{proposition}

\begin{proof}
(i) Let $j_0$ be the least index with $h_{j_0}\ne0$ and choose $u$ with $|u|
=1$ and $h_{j_0}(u)\ne0$. Each $\nabla h_j$ is homogeneous of degree $-d_j-1$,
so for $|w|\le\frac12$ and $\Lambda\ge1$, $|h_j(\Lambda u+w)-h_j(\Lambda u)|
\le|w|\sup_{|u'|\ge\Lambda/2}|\nabla h_j(u')|\le|w|(\Lambda/2)^{-d_j-1}
\sup_{|u'|=1}|\nabla h_j(u')|$. Put $y=-u$, so $|y|_\infty\ge\frac12$, take
$f,g\ge0$ in $C_c(\R^4)$ with $\int f=\int g=1$, $f$ supported in the ball of
radius $\frac14$ about $0$ and $g$ in that about $2y$. Then $(f,g)$ is
separated, with $\dist_\infty\ge\frac12$, and for $\lambda\ge0$ the pair
$(f,g(\cdot-\lambda y))$ is separated; for $x\in\operatorname{supp}f$ and
$x'-\lambda y\in\operatorname{supp}g$, $x-x'=\Lambda u+w$ with $\Lambda=
2+\lambda$ and $|w|\le\frac12$. Hence, if $S_2$ were given by $T$,
\[
\begin{split}
S_2\bigl(f,g(\cdot-\lambda y)\bigr)&=T(\Lambda u)+O(\Lambda^{-d_{j_0}-1})\\
&=h_{j_0}(u)\Lambda^{-d_{j_0}}+o(\Lambda^{-d_{j_0}}),
\end{split}
\]
which is at least $\frac12|h_{j_0}(u)|\Lambda^{-d_{j_0}}$ in absolute value
for large $\lambda$. On the other hand Theorem~\ref{thm:limitEsep}(v) with
$n=2$, $I=\{1\}$, $J=\{2\}$ and $S_1=0$ gives $|S_2(f,g(\cdot-\lambda y))|\le
Ke^{-m_*\lambda|y|_\infty}$ for large $\lambda$. The two are incompatible.

(ii) The two-point function of a Wick monomial ${:}\prod_{i=1}^j\partial^{
\alpha_i}\varphi{:}$ is, by Isserlis' formula and the orthogonality of Wick
monomials of different degrees~\cite{Jan97}, a sum over permutations $\sigma$
of $\pm\prod_i\partial^{\alpha_i}\partial^{\beta_{\sigma(i)}}G_0(x)$, for
monomials ${:}\prod_i\partial^{\alpha_i}\varphi{:}$ and ${:}\prod_i
\partial^{\beta_i}\varphi{:}$ of the same degree $j$ \textup{(}the two-point
function vanishes for different degrees\textup{)}, each term homogeneous of
degree $-2j-\sum_i(|\alpha_i|+|\beta_i|)<0$ and $C^\infty$ off $0$;
for $F_{\mu\nu}$ the covariance displayed follows from $\ip{A_\mu(x)A_\nu(0)}
=\delta_{\mu\nu}G_0(x)$ (Feynman gauge; the covariance of $F$ does not depend
on the gauge) by differentiating, and Wick monomials in the $F$ are of the
same type. So the two-point function of $\psi$ is a finite sum of homogeneous
terms of negative degrees, hence of the form $T$ of (i) after collecting
degrees, and nonzero by hypothesis; (i) applies. For the examples,
$\partial_\mu\partial_\nu|x|^{-2}=-2\delta_{\mu\nu}|x|^{-4}+8x_\mu x_\nu|x|^{-6}$
gives $\ip{F_{01}(x)F_{01}(0)}=-(\partial_0^2+\partial_1^2)G_0(x)$ as displayed,
and $\ip{\psi(x)\psi(0)}=2\sum_{\mu<\nu}\sum_{\rho<\sigma}\ip{F_{\mu\nu}(x)
F_{\rho\sigma}(0)}^2\ge2\ip{F_{01}(x)F_{01}(0)}^2$. For independent copies
the two-point function of a Wick polynomial is again of the form $T$.
\end{proof}

\begin{remark}\label{rN5:rem:massive}
The proof of Proposition~\ref{rN5:nomassless} uses only the exponential
clustering, and it cannot exclude Wick polynomials of \emph{massive} free
fields, which cluster exponentially; the massive free scalar field (the case
$\lambda=1$ of the construction above) and the massive free vector (Proca)
field satisfy the cube condition uniformly at fixed physical cube size
(Theorems~\ref{thm:ffscalar} and~\ref{thm:ffproca}). The short-distance
singularities treated in Propositions~\ref{prop:growth}
and~\ref{prop:limitEsing}, where (M) and (AF) enter, are a separate matter;
Proposition~\ref{rN5:nomassless} concerns the long-distance behaviour, where
only the gap enters.
\end{remark}

\subsubsection{Gaussian gauge fields and the string tension}

Let $A=(A_e)$ be a centred Gaussian field indexed by the oriented edges
$e$ of $\Z^4$, with $A_{-e}=-A_e$ and covariance $C(e,e')=\E A_eA_{e'}$. For a
closed lattice path $\gamma$ let $A(\gamma)=\sum_{e\in\gamma}A_e$, the edges
taken with the orientation of $\gamma$, and for $q\in\R$ let $W_q(\gamma)=\E
e^{iqA(\gamma)}=\exp(-\frac12q^2\Var A(\gamma))$. A \emph{lattice line} is
the set of edges of $\Z^4$ on a straight line parallel to a coordinate axis,
and
\[
K_\ell(C)=\sup_e\sup_\ell\sum_{e'\in\ell}|C(e,e')| ,
\]
the suprema over all edges $e$ and all lattice lines $\ell$. For integers
$R,T\ge1$ let $\gamma_{R,T}$ be the boundary of an $R\times T$ rectangle in a
coordinate plane, of perimeter $|\gamma_{R,T}|=2(R+T)$.

\begin{lemma}[perimeter law of Gaussian gauge fields]\label{rN5:gaussloop}
\textup{(i)} If $K_\ell(C)<\infty$, then for every rectangle,
$\Var A(\gamma_{R,T})\le4K_\ell(C)|\gamma_{R,T}|$ and $W_q(\gamma_{R,T})\ge
e^{-2q^2K_\ell(C)|\gamma_{R,T}|}$; in particular $\lim_{R\to\infty}R^{-2}\log
W_q(\gamma_{R,R})=0$ for every $q$.

\textup{(ii)} $K_\ell(C)<\infty$ holds whenever $|C(e,e')|\le K(1+|x_e-x_{e'}|
_\infty)^{-1-\eta}$ with $\eta>0$, $x_e$ the initial site of $e$, with
$K_\ell(C)\le K(1+2\sum_{n\ge1}n^{-1-\eta})$.

\textup{(iii)} Let $F$ be the infinite-volume form on $\Z^4$ of the free
Maxwell field $\gamma_\beta$ of Theorem~\textup{\ref{rD:maxwell}}: the
centred Gaussian field on the $2$-cochains with $\E\ip cF\ip{c'}F=\beta^{-1}
\ip{\Pi c}{\Pi c'}$ for finitely supported $c,c'$, $\Pi$ the orthogonal
projection of $\ell^2$ onto the closure of $B=dC^1$, and with Wilson loops
$\E e^{iq\ip{1_\Sigma}F}$ for a finite surface $\Sigma$ of oriented
plaquettes with boundary $\gamma$. This is the law of $F=dA$ for the
Feynman-gauge field $A$ with $C(e,e')=\beta^{-1}\delta_{\operatorname{dir}e,
\operatorname{dir}e'}G_0(x_e-x_{e'})$ \textup{(}$e,e'$ in the positive
orientation\textup{)}, $G_0=C_{0,(1,1,1,1)}$ of~\eqref{rN5:eq:green}, and
$\ip{1_\Sigma}{dA}=A(\gamma)$; so the Wilson loops of the Maxwell field are
the $W_q(\gamma)$ of this $A$, and $K_\ell(C)\le\beta^{-1}(K_0+\pi^2K_2/3)$
with $K_0,K_2$ of Lemma~\textup{\ref{rN5:green}(iii)} at $c=(1,1,1,1)$.
Hence the Maxwell field of $\U(1)^{\dim G}$ has perimeter law and string
tension zero at every $\beta$, and so does every Gaussian gauge field as in
\textup{(ii)}.
\end{lemma}

\begin{proof}
(i) We have
\[
\Var A(\gamma)=\sum_{e,e'\in\gamma}\pm C(e,e')\le\sum_{e\in\gamma}\sum_{s}
\sum_{e'\in s}|C(e,e')|,
\]
the middle sum over the four sides $s$ of the rectangle, each contained in a lattice line, so the right side is at most
$4K_\ell(C)|\gamma|$. (ii) A lattice line through the sites $x_0+ne$, $n\in\Z$,
contributes at most $K\sum_n(1+|n|)^{-1-\eta}$ up to a shift, bounded as
stated. (iii) Put $|\hat p|^2=\sum_\nu2(1-\cos p_\nu)$. The field $A$
exists, since $G_0$ is positive definite, being the Fourier transform of
$|\hat p|^{-2}\in L^1(T^4)$. In Fourier variables the operators $d$,
$\delta=d^{\mathsf T}$ and $-\Delta=d\delta+\delta d$ act on cochains as
matrices $\hat d(p)$, $\hat\delta(p)=\hat d(p)^*$ and $|\hat p|^2\mathbf 1$,
with $\hat d\hat d=0$, and $G_0$ acts on $1$-cochains as $|\hat p|^{-2}\mathbf
1$. For finitely supported $2$-cochains $b,b'$,
\[
\begin{split}
\E\ip b{dA}\ip{b'}{dA}&=\beta^{-1}\ip{\delta b}{G_0\delta b'}\\
&=\beta^{-1}\int_{T^4}\overline{\hat b}\cdot\hat P\hat b'\,
\frac{d^4p}{(2\pi)^4},
\end{split}
\]
where $\hat P(p)=\hat d\hat\delta/|\hat p|^2$. For $p\ne0$ the matrix $\hat
P$ is Hermitian, $\hat P\hat d=\hat d$ because
\[
\hat d\hat\delta\hat d=\hat d(|\hat p|^2-\hat\delta\hat d)=|\hat p|^2\hat d,
\]
and $\hat P$ vanishes on $\ker\hat\delta$, so $\hat P^2=\hat P$ and
the multiplier $P$ is an orthogonal projection on $\ell^2$, equal to $1$ on
$B$ and to $0$ on $\ker\delta=B^\perp$; hence $P=\Pi$, and $dA$ has the
covariance of $F$. The identity $\ip{1_\Sigma}{dA}=\ip{\delta1_\Sigma}A=
A(\partial\Sigma)$ is the definition of $d$. The bound on $K_\ell$ is (ii)
with Lemma~\ref{rN5:green}(iii) at $M=0$: $|G_0(z)|\le\min(K_0,K_2
|z|_\infty^{-2})$, and $\sum_{n\ne0}n^{-2}=\pi^2/3$.
\end{proof}

For the $\SU(N)$ Wilson states $\omega_k$ along a tuning $(\beta_k,a_k)$ put
$W_k(R,T)=\omega_k(N^{-1}\ReTr U(\gamma_{R,T}))$ and define the lattice string
tension $\sigma_k=\liminf_{R\to\infty}(-R^{-2}\log|W_k(R,R)|)\in[0,\infty]$.
\begin{enumerate}[label=\textup{(ST)},leftmargin=*]
\item \emph{Bounded physical string tension.} There are $0<s_-\le s_+<\infty$
and $k_0$ with $s_-a_k^2\le\sigma_k\le s_+a_k^2$ for all $k\ge k_0$.
\end{enumerate}

\begin{theorem}[a string tension excludes Gaussian gauge fields]\label{rN5:stgauss}
\textup{(i)} If the lower bound of \textup{(ST)} holds, then for every $k\ge
k_0$ there is no centred Gaussian gauge field $A$ with $K_\ell(C)<\infty$ and
no $q\in\R$ with $|W_k(R,R)|\ge W_q(\gamma_{R,R})$ for all large $R$; in
particular the rectangular Wilson loops of $\omega_k$ are not those of the
Maxwell field of $\U(1)^{\dim G}$ at any $\beta$, nor those of any Gaussian
gauge field with covariance decaying as in Lemma~\textup{\ref{rN5:gaussloop}(ii)},
and this holds at every $k$, not only in the limit.

\textup{(ii)} The cube condition and exponential clustering of the
gauge-invariant field do not by themselves give a positive string tension:
the Proca field $A$ on the links of $\Z^4$ with precision $\delta d+\mu^2$,
$\mu>0$, satisfies $\mathrm{BP}(n,\varpi)$ for gauge-invariant functions on
tori uniformly at fixed $n\mu$ \textup{(}Theorem~\textup{\ref{thm:ffproca}(ii)}
\textup{)}, its field strength $F=dA$ has the exponentially decaying covariance
$\E\ip bF\ip{b'}F=\ip b{d(-\Delta+\mu^2)^{-1}\delta b'}$, and
$\lim_RR^{-2}\log W_q(\gamma_{R,R})=0$ for every $q$.

\textup{(iii)} Suppose that along the tuning the strong-coupling area law
holds uniformly: $|W_k(R,T)|\le Ce^{-\sigma_0RT}$ for all $k$, $R$, $T$, with
$C<\infty$ and $\sigma_0>0$, as Osterwalder and Seiler proved for
$\beta_{W,k}$ in a range $[0,\beta_0]$~\cite{OSei78,Sei82}. Then the upper
bound of \textup{(ST)} fails, $\sigma_k/a_k^2\to\infty$, and every
rectangular Wilson loop of fixed physical size tends to zero, in accordance
with the triviality of the strong-coupling limits in
Corollary~\textup{\ref{cor:limitEstrong}}.
\end{theorem}

\begin{proof}
(i) If $|W_k(R,R)|\ge W_q(\gamma_{R,R})$ for large $R$,
Lemma~\ref{rN5:gaussloop}(i) gives
\[
\sigma_k\le\liminf_R\frac{-\log W_q(\gamma_{R,R})}{R^2}\le\liminf_R
\frac{8q^2K_\ell(C)}R=0<s_-a_k^2.
\] (ii) The cube condition is Theorem~\ref{thm:ffproca}(ii),
the Wilson loop variable $A(\gamma)$ being gauge invariant. On
$1$-cochains $(\delta d+\mu^2)^{-1}=(-\Delta+\mu^2)^{-1}P_T+\mu^{-2}P_L$ with
$P_L=d\delta(-\Delta)^{-1}$ the projection onto the exact cochains and $P_T=1-
P_L$ the projection onto $\ker\delta$, since $\delta d=-\Delta-d\delta$ vanishes
on exact cochains and equals $-\Delta$ on $\ker\delta$. For a closed path
$\gamma$ the indicator cochain $1_\gamma$ (signed by the orientation) lies in
$\ker\delta$, so
\[
\begin{split}
\Var A(\gamma)&=\ip{1_\gamma}{(-\Delta+\mu^2)^{-1}1_\gamma}\\
&=\sum_{e,e'\in\gamma}\pm\delta_{\operatorname{dir}e,\operatorname{dir}e'}
C_{\mu,(1,1,1,1)}(x_e-x_{e'}).
\end{split}
\]
Lemma~\ref{rN5:green}(iii), valid for every $\mu>0$, gives
$|C_{\mu,(1,1,1,1)}(z)|\le\min(K_0,K_2|z|_\infty^{-2})$, and the proof of
Lemma~\ref{rN5:gaussloop}(i), (ii) gives
$\Var A(\gamma_{R,R})\le16(K_0+\pi^2K_2/3)R$ and the claim on $W_q$.
Since $P_L\delta=0$ on $2$-cochains and $d$ commutes with $\Delta$, $d(\delta d+
\mu^2)^{-1}\delta=d(-\Delta+\mu^2)^{-1}\delta$, which is the covariance of $F$
as stated; the kernel $C_{\mu,(1,1,1,1)}$ decays exponentially by the same
lemma, applied with $\min(\mu,1)$ in the exponent $\varepsilon(0)$ (the
absolute value of the integrand of~\eqref{rN5:eq:3d} decreases in $M$). (iii) The bound gives
$\sigma_k\ge\sigma_0$, so $\sigma_k/a_k^2\ge\sigma_0a_k^{-2}\to\infty$, and a
loop of physical sides $\ell_1,\ell_2$ has $|W_k|\le Ce^{-\sigma_0\ell_1\ell_2/
a_k^2}\to0$.
\end{proof}

\begin{remark}[status]\label{rN5:rem:st}
(ST) along a continuum tuning is the confinement problem; it is not proved
here, nor is any area law for $\SU(N)$ at the weak couplings of a continuum
tuning. Theorem~\ref{rN5:stgauss} fixes the logical position of such an
input: the cube condition and the gap exclude the massless free theories
(Proposition~\ref{rN5:nomassless}) but not the massive ones, while a
positive physical string tension excludes every Gaussian gauge field with
summable covariance; the compact $\U(1)$ theory at weak coupling has a
perimeter law as well, by Guth~\cite{Gut80} and Fr\"ohlich and
Spencer~\cite{FS82}. The two inputs that remain for (E) beyond
Section~\ref{ssec:rEmissing} are thus dynamical: (W$_{\mathrm s}$), or
(SB$_n$) and (RI$_n$), for rotation invariance, and a nontriviality input
such as (ST), or a lower bound on a connected correlation of $\Phi$ that no
free field reproduces. By Theorem~\ref{rN5:nogo}, Lemma~\ref{rN5:bp} and
Remark~\ref{rN5:rem:nogo}, neither can be derived from reflection
positivity, the cube condition, the gap, clustering, uniqueness and
regularity alone, since fields at the sites have all these properties and
violate (W$_{\mathrm s}$), respectively are Gaussian; these fields are not
lattice gauge theories.
\end{remark}
